% ===== PRÉAMBULE CANONIQUE — à coller VERBATIM en tête de CHAQUE .tex =====
\documentclass[11pt,a4paper]{article}
\usepackage[utf8]{inputenc}\usepackage[T1]{fontenc}\usepackage{lmodern}
\usepackage[french]{babel}
\usepackage{amsmath,amssymb,mathtools}\usepackage{icomma}
\usepackage[version=4]{mhchem}          % \ce{...} : équations & formules
\usepackage{siunitx}\sisetup{locale=FR} % \qty{}{}, \si{}, \ang{}
\usepackage{chemfig}                    % \chemfig{...} : structures développées
\usepackage{xcolor}\usepackage{colortbl}
\usepackage{tikz}\usepackage{pgfplots}\pgfplotsset{compat=1.18}
\usepackage[most]{tcolorbox}\usepackage{enumitem}\usepackage{array,booktabs}
\usepackage[a4paper,margin=2.2cm]{geometry}
\usetikzlibrary{arrows.meta,calc,angles,quotes,positioning,patterns,backgrounds,shapes.geometric}
\definecolor{primary}{RGB}{222,49,99}\definecolor{primarydark}{RGB}{150,22,55}
\definecolor{defcol}{RGB}{0,114,178}\definecolor{methcol}{RGB}{52,92,148}
\definecolor{excol}{RGB}{0,135,90}\definecolor{attncol}{RGB}{200,40,55}
\tcbset{boxbase/.style={enhanced,breakable,boxrule=0.9pt,arc=2.5pt,
  left=3mm,right=3mm,top=2.3mm,bottom=2.3mm,fonttitle=\bfseries,coltitle=white}}
% --- boîtes TUTORIEL (noms EXACTS, ne pas renommer) ---
\newtcolorbox{cle}[1][]{boxbase,colback=defcol!6,colframe=defcol,
  title={\textsc{À retenir}\ifx\relax#1\relax\relax\else\textmd{ : \itshape #1}\fi}}
\newtcolorbox{methode}[1][]{boxbase,colback=methcol!7,colframe=methcol,sharp corners,
  title={\textsc{Méthode}\ifx\relax#1\relax\relax\else\textmd{ --- \itshape #1}\fi}}
\newtcolorbox{exemplebox}[1][]{boxbase,colback=excol!5,colframe=excol,
  title={\textsc{Exemple}\ifx\relax#1\relax\relax\else\textmd{ : \itshape #1}\fi}}
\newtcolorbox{piege}[1][]{boxbase,colback=attncol!6,colframe=attncol,
  title={\textsc{Piège}\ifx\relax#1\relax\relax\else\textmd{ : \itshape #1}\fi}}
\newtcolorbox{remarque}[1][]{boxbase,colback=black!4,colframe=black!45,
  title={\textsc{Remarque}\ifx\relax#1\relax\relax\else\textmd{ : \itshape #1}\fi}}
% --- boîtes EXERCICE / CORRIGÉ (noms EXACTS, auto-comptées) ---
\newtcolorbox[auto counter]{exo}[1][]{boxbase,colback=primary!4,colframe=primary,
  title={\textsc{Exercice}~\thetcbcounter\ifx\relax#1\relax\relax\else\textmd{ --- \itshape #1}\fi}}
\newtcolorbox[auto counter]{corrige}[1][]{boxbase,colback=excol!4,colframe=excol,
  title={\textsc{Corrigé}~\thetcbcounter\ifx\relax#1\relax\relax\else\textmd{ --- \itshape #1}\fi}}
\newcommand{\R}{\mathbb{R}}\newcommand{\N}{\mathbb{N}}\newcommand{\Z}{\mathbb{Z}}\newcommand{\ds}{\displaystyle}
% =========================================================================
\begin{document}
% THEME_TITLE: Le produit de solubilité Kps
\sisetup{per-mode=symbol}
\section*{\color{primarydark}Thème 2 — Le produit de solubilité $K_{\mathrm{ps}}$}

Le \textbf{produit de solubilité} $K_{\mathrm{ps}}$ est simplement la \emph{constante d'équilibre} de
la réaction de dissolution. Ce thème apprend à l'\textbf{écrire} correctement à partir de l'équation de
dissolution, puis à le \textbf{calculer} à partir des concentrations des ions. C'est une mécanique très
mécanique : bien poser l'équation, et le reste suit.

\begin{cle}[l'expression du produit de solubilité]
Pour un sel $\mathrm{M}_p\mathrm{X}_q$ qui se dissout selon
$\ce{M_{p}X_{q}(s) <=> p M^{q+}(aq) + q X^{p-}(aq)}$, le produit de solubilité est le
\textbf{produit des concentrations des ions}, chacune \textbf{élevée à son coefficient} :
\[ \boxed{\,K_{\mathrm{ps}} = [\ce{M^{q+}}]^{\,p}\,[\ce{X^{p-}}]^{\,q}\,} \]
\textbf{Règle d'or :} \emph{l'exposant de chaque ion dans $K_{\mathrm{ps}}$ est son coefficient dans
l'équation de dissolution.} Le solide, lui, n'apparaît pas.
\end{cle}

\begin{cle}[pourquoi le solide n'apparaît pas : activité $=1$]
L'\textbf{activité} d'un solide pur (ou du solvant eau) vaut \textbf{1} par convention : sa
« concentration efficace » ne dépend pas de la quantité présente. On l'« efface » donc de l'expression
de $K_{\mathrm{ps}}$ — seuls comptent les ions \emph{dissous}.
\end{cle}

\begin{methode}[écrire $K_{\mathrm{ps}}$ en deux temps]
\begin{enumerate}[leftmargin=1.7em,topsep=2pt,itemsep=1.5pt]
  \item \textbf{Écrire l'équation de dissolution} avec les bons coefficients $p$ et $q$.
  \item \textbf{Recopier} chaque concentration ionique en l'élevant à la puissance égale à son
        coefficient, et multiplier le tout.
\end{enumerate}
\end{methode}

\begin{exemplebox}[trois expressions typiques]
\begin{itemize}[leftmargin=1.6em,topsep=2pt,itemsep=2pt]
  \item $\ce{AgCl(s) <=> Ag^{+} + Cl^{-}}$ \ $\Rightarrow$ \ $K_{\mathrm{ps}} = [\ce{Ag+}]\,[\ce{Cl-}]$.
  \item $\ce{Fe(OH)3(s) <=> Fe^{3+} + 3 OH^{-}}$ \ $\Rightarrow$ \
        $K_{\mathrm{ps}} = [\ce{Fe^{3+}}]\,[\ce{OH-}]^{3}$.
  \item $\ce{Ag3PO4(s) <=> 3 Ag^{+} + PO4^{3-}}$ \ $\Rightarrow$ \
        $K_{\mathrm{ps}} = [\ce{Ag+}]^{3}\,[\ce{PO4^{3-}}]$.
\end{itemize}
\end{exemplebox}

\begin{piege}[coefficient = exposant, pas facteur]
Le « 3 » de \ce{Ag3PO4} devient un \textbf{exposant}, jamais un facteur multiplicatif :
\[ K_{\mathrm{ps}} = [\ce{Ag+}]^{3}[\ce{PO4^{3-}}] \quad\text{(correct)}, \qquad
   K_{\mathrm{ps}} \neq 3\,[\ce{Ag+}][\ce{PO4^{3-}}] \quad\text{(faux !)}. \]
\end{piege}

\begin{piege}[$K_{\mathrm{ps}}$ est SANS UNITÉ]
C'est une constante d'équilibre : on écrit $K_{\mathrm{ps}}(\ce{AgCl})=\num{1.8e-10}$ \emph{sans unité}.
Ne jamais lui coller des \si{\mole\per\litre} ni des « $\mathrm{M}^5$ ». En revanche la
\textbf{solubilité} $s$, elle, s'exprime bien en \si{\mole\per\litre}.
\end{piege}

\begin{cle}[cas particulier $MX_n$ (formule du livre)]
Pour un sel $\mathrm{MX}_n$ (un cation, $n$ anions), comme \ce{AgCl} ($n=1$), \ce{PbBr2} ($n=2$),
\ce{Fe(OH)3} ($n=3$) :
\[ K_{\mathrm{ps}}(\mathrm{MX}_n) = [\ce{M^{n+}}]\,[\ce{X^{-}}]^{\,n}. \]
\end{cle}

\begin{remarque}[calcul numérique]
Une fois l'expression écrite, calculer $K_{\mathrm{ps}}$ ne demande que de remplacer chaque
concentration \emph{à saturation} par sa valeur et de multiplier. Le résultat est un nombre pur (souvent
très petit, $10^{-10}$ ou moins).
\end{remarque}
\end{document}
